% POSTSCRIPT VERSION --- AFFECTS GRAPHICS AND OUTPUT ROUTINE 
% Copyright Michael J. Ferguson, INRS-Telecommunications
% All rights reserved. 
% Version Aug. 1991
% ---- TeX 3 version ----

%   Those items that are modified for different styles are indicated in 
%   each section. 


\input plain_nh 

% for TeX 3.+, hyphenation tables can be inserted in a format file only
% once. This means that both plain.tex rather than an old plain.fmt file
% must be used to create a new format file. plain_nh.tex is plain.tex with
% No Hyphenation (NH). 

% This is a cm font version ... now. 

% COMMENT: INRSTeX has been designed to allow for considerable variation
%   in the style of a document. This ranges from default spacing, fonts, 
%   formats of section heads, and headers and footers. The changes for
%   a style involves the changes of the defaults in "inrsdef.tex", document
%   "am" fonts in "inrsfont.tex" or "cm" fonts in inrsfncm.tex. 
%   The definitions of header and footer text is 
%   generally found in the various "xxsty.tex" files. These are incomplete
%   in the sense that only those changes to the defaults in "inrsdef.tex"
%   are actually done. Two examples are given, \paperstyle in "papersty.tex"
%   and \bookstyle in "booksty.tex". 
%

%   Those items that are modified for different styles are indicated in 
%   each section. 

\catcode`@= 11 % need to use some plain macros

%======= Additional active Characters ===========
% ALL ACTIVE CHARACTERS MUST BE DECLARED AND DEFINED AT THE VERY
% BEGINNING OR THERE IS A DANGER OF INCORRECT TOKENIZATION 
% ======== THESE ARE USED IN THE TABLE MACROS =========
\message{<< Additional Active Characters >>}
\catcode`\|= \active 
\def|{\ifmmode \vert\else \char`\|\fi} % effectively undoes activeness
\def\q@m{\string"}
\catcode`\"=\active \def"{\writeterm{<<< use `` or '' instead of \q@m >>>}{
                         }{}\char`\"}
%=== these are redefined inside the table macros =====


% ========= INRSTeX Utilities =========
% These are generally needed for good health. They include 
%    Houskeeping Macros 
%    Single/Multicolumn Page Eject and Insert Control 
%    Miscellaneous Paragraph Shape Macros 
%    Year/Month/Day 
%    New Boxit 
%    Number Format Conversion Forms 


%**************************
\input inrsutil.tex
%**************************


% ======== Extended Character Set(s) for Special Terminals =======
% The proper catcodes for the extended characters need to go in before the
% patterns are loaded. 

% *************************
   \ifundefined{charsubdef} 
           \message{ << Loading Extended Character Set -- TeX 3.+ >>}
            \input extdef3
          \else
            \message{ << Loading Extended Character Set -- MLTeX 3.+ >>} 
             \input extdef
           \fi
           \message { << Setting Parameters for TeX 3.+ >> }
           \emergencystretch =10pt 
           \lefthyphenmin=2
           \righthyphenmin=3
           \newcount\dischyph  % to effectively disable dischyph in mltex 
% *************************




% ======= Input Hyphenation patterns =======
% hyphenation for all languages are in the same file or accessed from it
% TeX 3 will not allow \patterns to be reopened once it has been closed.

         \input masthyph 


% ====== Multiple Language Support ========
% This may be done without bilingual TeX but is better done with it. 
% These macros define an \englishversion and a \versionfrancaise. 
% INRSTeX messages come in either English or French depending 
% upon the value of \language. If \language is not yet defined, a 
% \newcount is used to define it. 
% THIS MUST NOT BE CALLED BEFORE THE HYPHENATION PATTERNS ARE LOADED. 

%**************************
\input multlang.tex
%**************************


% ======  Basic Line Spacing =========
% The basic line spacing in a font is assumed to be scaled according to the 
% ex height of the font. This is perhaps not the aesthetically best but
% does allow consistent line spacing commands. The four basic commands are
% \normalbaselines (redefined from plain), 
% \normalbaselineskipsize, \spacing,  and \setstrut.

% ------- Defaults ... inrsdef.tex -----
% \def\normalbaselineskipsize{2.8ex \normalbaselineskipglue}
% \def\normalbaselineskipglue{plus .07ex minus .07ex}
% \def\normallineskipform{\lineskip 1pt \lineskiplimit 0pt}


%**************************
\input linespac.tex
%**************************



% ===== additional FONTS ==========
% The new fonts and font families are read in from the file "inrsfont.tex"
% These are read in by initex.
% This font list should define the font families. At the very least it must
% define fonts for the logical font names used in INRSTEX. These are
  %\notefont {usually \eighttt} ... foor version and margin notes.
 %\footerfont 
 %\headerfont
 %\footnotefont {should be a font family}
  %\smallerfont  {should be a font family}
  %\cheadfont 
  %\sheadfont
  %\ssheadfont
  %\sssheadfont
  %\dsssheadfont
  %\captionnumfont
  %\captiontextfont
  %\eqnumfont
  %\bigfont 
  %\biggfont 
  %\bigggfont 
  %\titlefont 

% The various \..head.. fonts are used in the \shead etc 
% The \footnotefont is used in footnotes and is a family 
% \titlefont is used in \title
% The other fonts are for convenience. 

% INRSTeX has the notion of a Document font. This allows for local changes in
% the font without messing up the sizes of the sectionheads, headers, footers
% etc. A font change inside a \documentstyle will cause the document font 
% to change.

%**************************
%\input inrsfont.tex %am fonts 
\input inrsfncm.tex % cm fonts 
%**************************



% ====== Marginnotes, Final and Draft Version =======
% INRSTeX's output routine supports two main versions of a document, a 
% "final version" and a "draft version". The "draft version" collects \refs
% and \cites in a margin box at the right of a page. In addition it puts on
% black marks for overfull lines and "Draft Version: <date> " at the
% bottom of the page.  

% If plain.tex output is to be used, "margform.tex" by "margform.pla"
% This will destroy all indications of \finalversion \draftversion

%**************************
\input margform.tex
%**************************




% =========== Document Styles, Page Numbering and Page Sizes ======
% These are the macros for setting page sizes and the various flags 
% the header/footer styles. This is only used with "inrsout.tex"

% ---- Defaults  ... set in  xxxsty.tex or inrsdef.tex  -------
%    Examples from "booksty.tex"
%       \vheadersize 
%       \vfootersize 
%       \headertext
%       \specialheadertext
%       \footertext  
%       \specialfootertext
%       \hpapersize 
%       \vpapersize 
%       \vouterpagesize 
%       \houterpagesize 
%       \topmargin 
%       \leftmargin



%**************************
\input pagsty.tex
%**************************




%======== Multicolumn format ===========
% These macros invoke multicolumnformat including balancing columns. 

% ---- Defaults ... set in inrsdef.tex 
%    \numcolumns 
%    \firstcolumnoffset 
%    \intercolumnsep 


%**************************
\input multicol.tex
%**************************



% ======== INRSTeX Single/Multi Column Output Routine ========
% This is the entire output routine for INRSTeX. It does both single and
% Multicolumn since there is a rather large repetition to do otherwise. 

%**************************
\input out_ps.tex
%**************************



%========= Autonumbering, Tagging and Autoreferencing =========

% These macros give the facility to symbolically refer to sections, equations,
% tables, and references and to automatically number them.
% Autonumber, automatically numbers sections, equations, ....
% Autoreferencing, causes \ref{tag} to drop the tag-value in the text.
% Autonumber and autoreference both true will cause tags to be generated. 
% Although it is possible to have only autonumbering, referring to 
% sections or other parts is quite difficult without referencing 

% ---- Defaults ... set in inrsdef.tex -----
%     \autoreferencetrue
%     \autonumberfalse  ... default when created
%     \def\Prerefform{}  \def\Postrefform  ... in auto 

%**************************
\input auto.tex
%**************************


% ============ Tag Generation ==========
% These are the macros that generate the actual tags. They are used in 
% automatic equation numbering, section heads and any \auto...num{tag} form. 

% ------ Defaults .... set in inrsdef.tex or xxxsty.tex ---
%       \undeftagmessage  ... set in "english.tex" or "francais.tex"


%**************************
\input tag.tex
%**************************


% ========== Equation Tags ===========
% Since equation numbering plays a big part in scientific documentation
% a special set of tag generation forms were produced for this case. 

% ------ Defaults .... set in inrsfont.tex -----
%    \eqnumfont 


%**************************
\input eqtag.tex
%**************************




%   ===== Citation generation ==========
% \cite is used with a journal/paper reference.
% The forms allow both random input order of cite forms or
% fixed order. At the moment it supports only IEEE form, ie 
% a list in the order of therir occurrence in a paper. 
% A data base system using a mailbox would be much better. 



%**************************
\input cite.tex
%**************************



% ========== Section, Chapter Headings  ===========
% These macros have several parts,
%       -- Counter allocation and numbering styles for autonumber
%       -- The format of the heading. 
%       -- The writing of a list file for automatic table of contents.
%       -- The referencing in the text ... ie reference number styles.

% -------- Defaults ... inrsdef.tex, inrsfont.tex, xxxsty.tex ----
%   
%   \chnum=0   \def\chnumform{\the\chnum}  ... inrsdef.tex
%   \shnum=0   \def\shnumform{\the\shnum}  ... inrsdef.tex
%   \sshnum=0   \def\sshnumform{\the\sshnum}  ... inrsdef.tex
%   \ssshnum=0 \def\ssshnumform{\the\ssshnum}  ... inrsdef.tex
%   \def\chtagrefformat{\chnumform}  ... secthead.tex 
%   \def\chtagreplaceformat{\Chapter \chtagrefformat}  .. secthead.tex
%      \prsheadskip= 6ex plus 2ex minus 2ex  ... inrsdef.tex 
%      \posheadskip= 2ex   ... inrsdef.tex 
%      \prssheadskip= 5ex plus 1.8ex minus 1.8ex  ... inrsdef.tex 
%      \possheadskip= 2ex  ... inrsdef.tex 
%      \prsssheadskip= 3ex  plus 1ex minus 1ex  ... inrsdef.tex 
%      \posssheadskip= 2ex  ... inrsdef.tex 
%      \prdsssheadskip= 3ex  plus 1ex minus 1ex  ... inrsdef.tex 
%      \podsssheadskip= 2ex  ... inrsdef.tex 
%
%      \gensheadformat#1#2#3#4 ... secthead.tex ... for major format
%                                                      surgery.
%
%      \nochaptertrue  ... inrsdef.tex
%      \def\Chapter ... english.tex or francais.tex 
%      \cheadformat#1#2  ... secthead.tex 
%      \def\chapterstartform{} ... secthead.tex and xxxsty.tex


%**************************
\input secthead.tex
%**************************


%============ Preface/Prelude macros ===========
% this is the first of the environment macros. The problem, among others is
% that the headers, footers, pagenumbering, etc are different in preface of
% the document than elsewhere. 

% ---- Defaults ... prelhead.tex or xxxsty.tex 
%      \preludeheadformat#1  ... prelhead.tex 

%**************************
\input prelhead.tex
%**************************



% =========== Appendices  ==============
% Appendices are defined by using \shead invocations.
% Appendices are assumed to be associated with chapters, if 
% they exist  hence the first appendix in Chapter 2 is Appendix A.
% These macros modify the numberforms of \shead to produce Apendices. 
% The "appendix head" is an \shead with a different numberform. 

% ------- Defaults ------
%     \def\Appendix  ... english.tex or francais.tex 

%**************************
\input appendix.tex
%**************************



% ==========  List Files =========
% General macros for making, open and closing "list' files. These are files
% where the form output includes a page reference. They are named with a 
% three letter extension ... eg ".fig"  ... to the \jobname form. 

%**************************
\input writfile.tex
%**************************



% ======== Table of Contents, Figures, and Tables -- Examples ======
% The "list" file is written using \writelistfile and the \shtoc ...
% are used to create the actual table of contents, figures, tables. 
% Extensions to other forms should be obvious. 

% ----- Defaults ----
%      \righttocindent  ... inrsdef.tex 
%      \tocfill  ... tocform.tex  ... dot form for toc list.
%      tocfonts ... tocform.tex ... in actual \figtoc etc

%**************************
\input tocform.tex
%**************************





%======== List macros ==========
% there is only one kind of list in this system, namely on that is
% right justified in the particular box that the list appears
% extensions may be made to include a center

% -----  Defaults  ... inrsdef.tex -----
%    \prlistskip= 2ex plus 3pt minus 2pt 
%    \prsublistskip = 1ex plus 2pt minus 1pt
%    \prsubsublistskip = .5ex plus 1pt minus 1pt
%    \polistskip= 2ex plus 1pt minus 1pt
%    \posublistskip= 1ex plus 1pt minus 1pt
%    \posubsublistskip= .5ex plus 1pt minus 1pt

%    \listitemskip = 2ex plus 1pt minus 1pt
%    \sublistitemskip =1ex plus .75pt minus .75pt
%    \subsublistitemskip = .5ex plus .5pt minus .5pt

%    \listindent = 3em
%    \sublistindent = 6em
%    \subsublistindent = 9em
%    \listitemmarksize = 1.5em

%**************************
\input lists.tex
%**************************



%=========== Footnotes ---- similar to that in the  texbook ======

% ------ Defaults --------
%      \footnotefont  ... family in inrsfont.tex 
%      \skip\footins  ... footnote.tex 
%      \count\footins ... plain.tex 
%      \dimen\footins ... footnote.tex 
%      \footnoterule  ... inrsout.tex 


%**************************
\input footnote.tex
%**************************



%========= ABSTRACT/TITLE  =======
% These generate a centered abstract in the \smallerfont of a family 
% and define the \title macro 

% ------ Defaults ------
%     \absindent ... inrsdef.tex  abstract indentation
%     \ABSTRACT  ... english.tex or francais.tex 
%     \titlemessage ... english.tex of francais.tex 
%     \titlefont ... inrsfont.tex 
%     \titlefraction ... inrsdef.tex 
%     \titlespacing  ... inrsdef.tex 

%**************************
\input abstract.tex
%**************************


% ========== Captioned Table/Figure Forms and Inserts =========
% Insert Forms 

%   ----- Defaults -----
%    \captiontextfraction  ... inrsdef.tex 
%    \captionskip  ... inrsdef.tex 
%    \topcaptiontrue or \topcaptionfalse  ... inrsdef.tex 
%    \figurelabel  or \figlabel  ... english.tex or francais.tex 
%    \tablelabel   or \tbllabel  ... english.tex or francais.tex 
%    \captionnumfont ... inrsfont.tex
%    \captiontitlefont ... inrsfont.tex
%    \captionbodyfont ... inrsfont.tex
%    \captionbodyon <off> ... inrsdef.tex 
%    \silentfigurefalse 

%**************************
\input genins.tex
%**************************


% ============  Table Making Forms ============

%**************************
\input tables.tex
%**************************



%========= Verbatim or NoFill Style ========
% The verbatim form is not part of the INRSTeX format file. 
% The  following brings in the verbatim macros and redefines \beginttverbatim

\def\beginttverbatim{\input verbatim\relax}


%========== Special styles ===============
%----- Special Document Styles/Forms  -------
% INRSTeX learns of the existence of special document styles or forms
% from specform.tex

%**************************
\input specform.tex
%**************************



%======== Multiple part documents -- subdocuments =======
% documents may be broken down to the chapter or section level 

%**************************
\input subdoc.tex
%**************************





% ============ TeXGraphics Support =============
% TeXGraphics uses \specials, assumes POSTSCRIPT and assumes a driver
% equivalent to the INRS modified version of Beebe's dvialw. 

%**************************
\input graph_ps.tex
%**************************






% ======== Bugs/Modifications to Plain =====

%**************************
\input plainmod.tex
%**************************



% ====== INRSTeX defaults .... in addition to Plain ....
%**************************
\input inrsdef.tex
%**************************


%  =========== modifications for MSDOS usage ============

% for Unix systems replace the \input writdos by 

 \message{Input Unix write.file form} 
 \input writunix

% invoke inrstex with the shell script 
%  inrsunix.sh ... appropriately modified 


%\message{Input DOS write/file form}

%\input writdos


\catcode`@=12 % puts at  back to a non letter

\def\fmtname{INRSTeX-Bilingual-PostScript}\def\fmtversion{3.1} 

